\chapter{Ablaciones complementarias sobre DermapixelAI 1.0}
\label{cap:anexoj}

Este anexo recoge varios estudios complementarios sobre
DermapixelAI~1.0 que confirman, sin desplazarla, la línea
experimental principal del capítulo~\ref{cap:resultados}
(secciones~\ref{sec:dermapixel-results} y~\ref{sec:dermapixel-zs}):
la validación del procedimiento de etiquetado del dataset, tres
variantes de clasificación sobre \emph{embeddings} congelados
---cabezas alternativas, función de pérdida y \emph{Test-Time
Augmentation}--- y dos extensiones con codificadores generalistas y
arquitecturas de inferencia alternativas. Ninguno modifica las
conclusiones del cuerpo principal; se documentan aquí por
completitud y trazabilidad.


\section{Validación del procedimiento de etiquetado: subestudio
\emph{rosa\_verified}}
\label{sec:anexoj-rosa}

La sección~\ref{sec:dermapixel-results} resume el subestudio que
cuantifica si las etiquetas etiológicas derivadas de la ontología,
no revisadas individualmente por la dermatóloga, degradan el
rendimiento. La tabla~\ref{tab:dermapixel-rosa} reporta la
comparación completa, por sondeo lineal, entre el subconjunto
verificado caso a caso (\texttt{rosa\_verified=True}, $N = 919$;
train $= 744$, test $= 32$) y el dataset completo
(\texttt{all\_verified}, $N = 1\,062$; train $= 874$, test $= 36$),
sobre los tres niveles ontológicos.

\begin{table}[H]
\centering
\footnotesize
\setlength{\tabcolsep}{4pt}
\renewcommand{\arraystretch}{0.95}
\caption{Subestudio \emph{rosa\_verified}: comparación entre el
subconjunto verificado caso a caso por la dermatóloga
($N_{\mathrm{train}} = 744$, $N_{\mathrm{test}} = 32$) y el dataset
completo ($N_{\mathrm{train}} = 874$, $N_{\mathrm{test}} = 36$).
Entre corchetes, IC95\% por \emph{bootstrap} estratificado.}
\label{tab:dermapixel-rosa}
\begin{tabular}{lllrrr}
\hline
Modelo & Subconjunto & Nivel & Acc & BAcc & W-F1 \\
\hline
PanDerm Large & \emph{rosa\_true\_only} & L1
  & $0{,}719$ [$0{,}594$--$0{,}844$]
  & $0{,}747$ [$0{,}646$--$0{,}865$]
  & $0{,}700$ [$0{,}550$--$0{,}841$] \\
PanDerm Large & \emph{all\_verified}    & L1
  & $0{,}750$ [$0{,}639$--$0{,}861$]
  & $0{,}735$ [$0{,}637$--$0{,}835$]
  & $0{,}722$ [$0{,}596$--$0{,}845$] \\
PanDerm Large & \emph{rosa\_true\_only} & L2
  & $0{,}250$ [$0{,}188$--$0{,}344$]
  & $0{,}250$ [$0{,}203$--$0{,}313$]
  & $0{,}227$ [$0{,}150$--$0{,}304$] \\
PanDerm Large & \emph{all\_verified}    & L2
  & $0{,}250$ [$0{,}167$--$0{,}333$]
  & $0{,}265$ [$0{,}206$--$0{,}338$]
  & $0{,}223$ [$0{,}139$--$0{,}291$] \\
PanDerm Large & \emph{rosa\_true\_only} & L3
  & $0{,}167$ [$0{,}167$--$0{,}167$]
  & $0{,}176$ [$0{,}176$--$0{,}176$]
  & $0{,}125$ [$0{,}117$--$0{,}139$] \\
PanDerm Large & \emph{all\_verified}    & L3
  & $0{,}179$ [$0{,}143$--$0{,}214$]
  & $0{,}184$ [$0{,}158$--$0{,}211$]
  & $0{,}143$ [$0{,}100$--$0{,}184$] \\
PanDerm Base  & \emph{rosa\_true\_only} & L1
  & $0{,}625$ [$0{,}469$--$0{,}781$]
  & $0{,}523$ [$0{,}326$--$0{,}735$]
  & $0{,}618$ [$0{,}451$--$0{,}773$] \\
PanDerm Base  & \emph{all\_verified}    & L1
  & $0{,}639$ [$0{,}472$--$0{,}778$]
  & $0{,}570$ [$0{,}371$--$0{,}751$]
  & $0{,}624$ [$0{,}457$--$0{,}760$] \\
\hline
\end{tabular}
\end{table}


\section{Ranking L3 por prototipos coseno de Dermapixel R0}
\label{sec:anexoj-l3-prototipos}

La sección~\ref{sec:dermapixel-spanderm-v0} resume el ranking del
nivel L3 de Dermapixel R0 por prototipos coseno (media de las
representaciones de entrenamiento por clase). La
tabla~\ref{tab:dermapixel-spanderm-l3} reporta las cifras completas
sobre las $250$ clases efectivas del vocabulario de entrenamiento.

\begin{table}[H]
\centering
\footnotesize
\setlength{\tabcolsep}{4pt}
\renewcommand{\arraystretch}{0.95}
\caption{Ranking L3 por prototipos coseno de Dermapixel R0 sobre el
conjunto de test de DermapixelAI~1.0 ($250$ clases efectivas
del vocabulario de entrenamiento). Cifras: media $\pm$ desviación
estándar sobre tres semillas. La columna $\times$~azar reporta el
factor de mejora respecto a la asignación aleatoria uniforme
sobre $250$ clases ($1/250 = 0{,}004$).}
\label{tab:dermapixel-spanderm-l3}
\begin{tabular}{lrrr}
\hline
Métrica & Media                 & Std            & $\times$~azar \\
\hline
Acc@1   & $0{,}167$             & $0{,}000$      & $\times 42$ \\
Acc@3   & $0{,}259$             & $0{,}013$      & $\times 22$ \\
Acc@5   & $0{,}296$             & $0{,}026$      & $\times 15$ \\
BAcc    & $0{,}148$             & $0{,}000$      & --- \\
\hline
\end{tabular}
\end{table}

\section{Cabezas alternativas sobre embeddings congelados}

Sobre los mismos \emph{embeddings} de PanDerm Large utilizados en la
tabla~\ref{tab:dermapixel-lp} se evalúan tres cabezas adicionales:
$k$-vecinos más cercanos con distancia coseno ($k = 10$, valor
óptimo dentro de $k \in \{1, 5, 10\}$), perceptrón multicapa con
una capa oculta de $512$ unidades y \emph{dropout} $0{,}3$ (MLP
512), y regresión logística con \texttt{class\_weight='balanced'}
(LogReg \emph{balanced}). El protocolo de partición es el del split
fijo de la tabla~\ref{tab:dermapixel-lp}
($N_{\mathrm{train}} = 874$, $N_{\mathrm{test}} = 36$ en L1/L2,
$N_{\mathrm{test}} = 28$ en L3). La tabla~\ref{tab:dermapixel-heads}
recoge la AUROC \emph{one-vs-rest}, métrica discriminativa robusta
al desbalance, como criterio comparativo.

\begin{table}[H]
\centering
\footnotesize
\setlength{\tabcolsep}{4pt}
\renewcommand{\arraystretch}{0.95}
\caption{Cabezas alternativas sobre \emph{embeddings} congelados de
PanDerm Large (test fijo de DermapixelAI~1.0). AUROC
\emph{one-vs-rest} por nivel ontológico. LP: regresión logística
estándar (referencia tabla~\ref{tab:dermapixel-lp}); kNN $k=10$:
$k$-vecinos con distancia coseno; MLP 512: perceptrón multicapa
con $512$ unidades ocultas y \emph{dropout} $0{,}3$; LogReg
\emph{balanced}: regresión logística con \texttt{class\_weight =
balanced}. Mejor valor por nivel en negrita.}
\label{tab:dermapixel-heads}
\begin{tabular}{lrrr}
\hline
Cabeza                & L1 (4 cls)         & L2 (38 cls)        & L3 (224 cls) \\
\hline
LP (referencia)       & $0{,}796$          & $0{,}793$          & $0{,}813$ \\
kNN $k=10$            & $0{,}752$          & $0{,}698$          & $0{,}694$ \\
MLP 512               & $\mathbf{0{,}817}$ & $\mathbf{0{,}812}$ & $\mathbf{0{,}827}$ \\
LogReg \emph{balanced}& $0{,}797$          & $0{,}800$          & $0{,}810$ \\
\hline
\end{tabular}
\end{table}

El MLP de $512$ unidades alcanza la mejor AUROC sobre los tres
niveles, con incrementos modestos pero consistentes sobre la cabeza
lineal de referencia: $+2{,}1$\,pp en L1
($0{,}796 \to 0{,}817$), $+1{,}9$\,pp en L2
($0{,}793 \to 0{,}812$) y $+1{,}4$\,pp en L3
($0{,}813 \to 0{,}827$). La regresión logística con re-ponderación
de clases mantiene Acc@1 idéntico al LP estándar y mejora la BAcc
sobre L1 ($0{,}703$ frente a $0{,}735$ del LP estándar dentro de
los IC95\% solapados), sin diferencias relevantes sobre L2 y L3.
La cabeza $k$-NN queda por debajo de las cabezas paramétricas en
los tres niveles, lo que se interpreta como insuficiente densidad
local en el espacio de \emph{embeddings} para $N_{\mathrm{train}}
= 874$ frente a un vocabulario de $224$ clases L3. La elección de
la regresión logística estándar como cabeza primaria en la
tabla~\ref{tab:dermapixel-lp} se mantiene defendible: la mejora
absoluta del MLP no compensa la pérdida de interpretabilidad
lineal de los coeficientes ni la mayor sensibilidad a la
inicialización aleatoria, especialmente con el tamaño muestral del
test.

\section{Función de pérdida: Cross-Entropy frente a Focal}
\label{sec:dermapixel-focal}

La cola larga severa del dataset, particularmente acusada en el
nivel L3 ($77\%$ de las clases con $\leq 5$ ejemplos de
entrenamiento, $3{,}9$ muestras por clase de media en
\emph{train}, ver Anexo~\ref{cap:anexoc}), motiva contrastar la
pérdida de entropía cruzada estándar con \emph{Focal Loss} de
Lin et~al.\ 2017, diseñada explícitamente para mitigar el efecto
del desbalance reduciendo el peso de los ejemplos fáciles vía el
factor $(1-p_t)^\gamma$. Sobre los \emph{embeddings} congelados de
PanDerm Large del split fijo de la tabla~\ref{tab:dermapixel-lp}
se reentrena la cabeza lineal con \emph{Focal Loss} bajo dos
parámetros de focalización ($\gamma \in \{1{,}0,\,2{,}0,\,5{,}0\}$)
y dos regímenes de ponderación de clase
(\texttt{class\_weight = None} frente a
\texttt{class\_weight = balanced}), optimizada con LBFGS y
evaluada sobre el mismo conjunto de test que el sondeo lineal
estándar. La tabla~\ref{tab:dermapixel-focal} reporta la mejor
variante por nivel junto a las dos referencias previas (LP CE
estándar, LogReg \emph{balanced}) de la
tabla~\ref{tab:dermapixel-heads}.

\begin{table}[H]
\centering
\footnotesize
\setlength{\tabcolsep}{4pt}
\renewcommand{\arraystretch}{0.95}
\caption{Función de pérdida sobre \emph{embeddings} congelados de
PanDerm Large (test fijo de DermapixelAI~1.0,
$N_{\mathrm{test}} = 36$ en L1/L2 y $N_{\mathrm{test}} = 28$ en
L3). LP CE: regresión logística con entropía cruzada estándar
(referencia tabla~\ref{tab:dermapixel-lp}); LogReg
\emph{balanced}: re-ponderación inversa por frecuencia de clase
(tabla~\ref{tab:dermapixel-heads}); \emph{Focal} $\gamma$: pérdida
focal sin re-ponderación con parámetro de focalización indicado.
Mejor valor por columna y nivel en negrita.}
\label{tab:dermapixel-focal}
\begin{tabular}{llrrr}
\hline
Nivel        & Pérdida                                       & Acc@1              & BAcc               & AUROC \\
\hline
L1 (4 cls)   & LP CE (referencia)                            & $0{,}750$          & $0{,}735$          & $0{,}796$ \\
L1 (4 cls)   & LogReg \emph{balanced}                        & $0{,}694$          & $0{,}703$          & $0{,}797$ \\
L1 (4 cls)   & \emph{Focal} $\gamma = 2{,}0$                 & $\mathbf{0{,}778}$ & $\mathbf{0{,}768}$ & $\mathbf{0{,}800}$ \\
\hline
L2 (38 cls)  & LP CE (referencia)                            & $0{,}250$          & $0{,}265$          & $\mathbf{0{,}793}$ \\
L2 (38 cls)  & LogReg \emph{balanced}                        & $0{,}250$          & $0{,}265$          & $\mathbf{0{,}800}$ \\
L2 (38 cls)  & \emph{Focal} $\gamma = 2{,}0$                 & $\mathbf{0{,}278}$ & $\mathbf{0{,}324}$ & $0{,}778$ \\
L2 (38 cls)  & \emph{Focal} $\gamma = 5{,}0$                 & $\mathbf{0{,}278}$ & $\mathbf{0{,}324}$ & $\mathbf{0{,}793}$ \\
\hline
L3 (224 cls) & LP CE (referencia)                            & $\mathbf{0{,}143}$ & $\mathbf{0{,}132}$ & $\mathbf{0{,}813}$ \\
L3 (224 cls) & \emph{Focal} $\gamma = 2{,}0$                 & $\mathbf{0{,}143}$ & $\mathbf{0{,}132}$ & $0{,}771$ \\
\hline
\end{tabular}
\end{table}

\emph{Focal Loss} con $\gamma = 2{,}0$ y sin re-ponderación de
clase supera marginalmente a la entropía cruzada estándar sobre L1
($+3{,}3$\,pp BAcc) y L2 ($+5{,}9$\,pp BAcc), pero degrada la
AUROC L2 en $-1{,}5$\,pp ($0{,}793 \to 0{,}778$). Añadir
\texttt{class\_weight = balanced} no aporta, porque el factor
$(1-p_t)^\gamma$ ya modula el peso por dificultad. Sobre L3 ningún
ajuste de la pérdida mejora a la entropía cruzada (Acc@1 y BAcc
idénticos en $0{,}143$ y $0{,}132$; AUROC $-4{,}2$\,pp). Como en
el FT denso (sección~\ref{sec:dermapixel-ft-l1}), la cola larga
estructural de L3 no admite mitigación por la función de pérdida
sobre el dataset actual.

El patrón observado es análogo al documentado en el aumento en
tiempo de inferencia (TTA, sección~\ref{sec:dermapixel-tta}):
\emph{Focal Loss} desplaza el régimen Acc@1/BAcc-AUROC
favoreciendo la primera dimensión sobre la segunda, pero no
resuelve el cuello de botella estructural impuesto por la
cardinalidad y la cola larga del dataset. La elección del LP CE
estándar como cabeza primaria de la tabla~\ref{tab:dermapixel-lp}
se mantiene defendible: la mejora de \emph{Focal Loss} sobre L1 y
L2 está dentro de los IC95\% solapados y no compensa la pérdida
de interpretabilidad probabilística calibrada de la entropía
cruzada estándar.

\section{Aumento en tiempo de inferencia}
\label{sec:dermapixel-tta}

La sección~\ref{sec:ensemble-results} documenta la aplicación de
\emph{Test-Time Augmentation} (TTA) sobre el clasificador M1
(PanDerm Large \emph{fine-tuned}) del servidor de inferencia, con un
patrón de cinco augmentaciones deterministas (original, simetría
horizontal, simetría vertical, rotación de $90^\circ$ y rotación de
$270^\circ$, omitiendo la rotación de $180^\circ$). El presente
experimento replica el protocolo sobre la cabeza lineal entrenada
sobre los \emph{embeddings} de PanDerm Large del dataset
DermapixelAI~1.0: para cada imagen de test se extraen los cinco
\emph{embeddings} correspondientes a las cinco augmentaciones, cada
uno se evalúa con la cabeza lineal entrenada sobre el conjunto de
\emph{train} original y las cinco distribuciones de probabilidad se
promedian antes del \emph{argmax} final. La
tabla~\ref{tab:dermapixel-tta} compara el sondeo lineal sin TTA
(referencia tabla~\ref{tab:dermapixel-lp}) frente al sondeo lineal
con TTA sobre los tres niveles ontológicos.

\begin{table}[H]
\centering
\footnotesize
\setlength{\tabcolsep}{4pt}
\renewcommand{\arraystretch}{0.95}
\caption{Sondeo lineal con PanDerm Large sobre DermapixelAI~1.0:
referencia sin TTA frente a TTA con cinco augmentaciones
deterministas (original, simetría horizontal, simetría vertical,
rotación $90^\circ$, rotación $270^\circ$). Promedio de
probabilidades sobre las cinco augmentaciones antes del
\emph{argmax}. Mejor valor por par en negrita.}
\label{tab:dermapixel-tta}
\begin{tabular}{llrrrr}
\hline
Nivel & Modo & Acc@1 & Acc@3 & BAcc & AUROC \\
\hline
L1 (4 cls)   & LP                & $\mathbf{0{,}750}$ & $0{,}944$          & $\mathbf{0{,}735}$ & $\mathbf{0{,}796}$ \\
L1 (4 cls)   & LP $+$ TTA        & $0{,}694$          & $\mathbf{1{,}000}$ & $0{,}703$          & $0{,}779$ \\
L2 (38 cls)  & LP                & $\mathbf{0{,}250}$ & $0{,}278$          & $\mathbf{0{,}265}$ & $\mathbf{0{,}793}$ \\
L2 (38 cls)  & LP $+$ TTA        & $0{,}222$          & $\mathbf{0{,}500}$ & $0{,}250$          & $0{,}787$ \\
L3 (224 cls) & LP                & $\mathbf{0{,}179}$ & $0{,}214$          & $\mathbf{0{,}184}$ & $\mathbf{0{,}813}$ \\
L3 (224 cls) & LP $+$ TTA        & $0{,}179$          & $\mathbf{0{,}250}$ & $0{,}184$          & $0{,}807$ \\
\hline
\end{tabular}
\end{table}

El patrón es consistente sobre los tres niveles ontológicos: TTA
degrada Acc@1, BAcc y AUROC en magnitudes pequeñas ($-5{,}6$,
$-2{,}8$ y $0{,}0$\,pp Acc@1 sobre L1, L2 y L3 respectivamente;
$-3{,}2$, $-1{,}5$ y $0{,}0$\,pp BAcc; $-1{,}7$, $-0{,}6$ y
$-0{,}6$\,pp AUROC) y mejora sistemáticamente Acc@3
($+5{,}6$\,pp en L1 hasta alcanzar $1{,}000$, $+22{,}2$\,pp en
L2 y $+3{,}6$\,pp en L3). El comportamiento es coherente con la
observación reportada en sección~\ref{sec:ensemble-results} para
M1 sobre HAM10000, donde TTA produce $-2{,}6$\,pp BAcc y
$+4{,}3$\,pp \emph{recall} de melanoma. La interpretación es que el
promediado sobre cinco augmentaciones suaviza la distribución de
probabilidad y reduce la confianza en la clase \emph{argmax} pero
distribuye masa de probabilidad entre clases visualmente cercanas,
lo que penaliza la precisión \emph{top-1} y mejora el conjunto
\emph{top-3}. PanDerm Base presenta la misma dirección de efecto
(cifras completas en el material reproducible). La recomendación
operativa derivada es que TTA es preferible en escenarios de cribado
con presentación \emph{top-k} de diagnósticos diferenciales al
clínico (Acc@3 L2 $= 0{,}500$ frente a $0{,}278$ sin TTA), y el
sondeo lineal sin TTA es preferible para decisiones automatizadas
\emph{top-1}.


\section{Codificadores generalistas no médicos: DINOv2 y CLIP-L}
\label{sec:dermapixel-extra-encoders}

Se evalúan dos codificadores fuera del dominio biomédico bajo el
protocolo de sondeo lineal congelado descrito en la
sección~\ref{sec:dermapixel-results}, sobre el mismo
\emph{split} fijo ($N_{\mathrm{train}} = 874$,
$N_{\mathrm{test}} = 36$ en L1/L2, $N_{\mathrm{test}} = 28$ en
L3). DINOv2 ViT-L (\texttt{vit\_large\_patch14\_dinov2.lvd142m},
embedding $1\,024$-D, CLS) corresponde al modelo auto-supervisado
visual generalista de Meta sobre LVD-142M~\cite{dinov2}; CLIP-L-336
(OpenAI, ViT-L/14 con resolución $336$) corresponde al modelo
contrastivo texto-imagen generalista entrenado sobre $400$
millones de pares web~\cite{clip2021}, con embedding proyectado de
$768$-D. La cabeza es regresión logística L-BFGS con $C = 1{,}0$ y
$\mathrm{max\_iter} = 5\,000$, idéntica a la
sección~\ref{sec:dermapixel-results}. La
tabla~\ref{tab:dermapixel-extra-encoders} recoge los resultados.

\begin{table}[H]
\centering
\footnotesize
\setlength{\tabcolsep}{4pt}
\renewcommand{\arraystretch}{0.95}
\caption{Sondeo lineal sobre embeddings congelados de
codificadores generalistas no médicos en DermapixelAI~1.0
($N_{\mathrm{test}} = 36$ en L1/L2, $N_{\mathrm{test}} = 28$ en
L3). Referencia: PanDerm Large
(tabla~\ref{tab:dermapixel-lp}). Mejor valor por columna y nivel en
negrita.}
\label{tab:dermapixel-extra-encoders}
\begin{tabular}{lrrrrrr}
\hline
              & \multicolumn{2}{c}{L1 (4 cls)} & \multicolumn{2}{c}{L2 (38 cls)} & \multicolumn{2}{c}{L3 (224 cls)} \\
Encoder       & BAcc & AUROC & BAcc & AUROC & BAcc & AUROC \\
\hline
PanDerm Large & $\mathbf{0{,}735}$ & $\mathbf{0{,}796}$
              & $0{,}265$          & $0{,}793$
              & $0{,}184$          & $0{,}813$ \\
DINOv2 ViT-L  & $0{,}537$          & $0{,}650$
              & $0{,}250$          & $0{,}718$
              & $\mathbf{0{,}211}$ & $0{,}794$ \\
CLIP-L-336    & $0{,}556$          & $0{,}658$
              & $\mathbf{0{,}279}$ & $\mathbf{0{,}826}$
              & $0{,}184$          & $\mathbf{0{,}827}$ \\
\hline
\end{tabular}
\end{table}

CLIP-L-336 alcanza la mejor AUROC sobre L2 ($0{,}826$ frente a
$0{,}793$ de PanDerm Large, $+3{,}3$\,pp) y sobre L3
($0{,}827$ frente a $0{,}813$, $+1{,}4$\,pp) a pesar de no
incorporar material biomédico en su preentrenamiento. DINOv2 ViT-L,
auto-supervisado puro sin componente textual, gana la BAcc L3
($0{,}211$ frente a $0{,}184$ de PanDerm Large,
$+2{,}7$\,pp) pero pierde sistemáticamente sobre L1
($-19{,}8$\,pp BAcc, $-14{,}6$\,pp AUROC) y se mantiene por
debajo sobre L2 ($-1{,}5$\,pp BAcc, $-7{,}5$\,pp AUROC).
PanDerm Large conserva la primacía sobre L1 en BAcc y AUROC.
Ningún codificador domina simultáneamente en todas las métricas y
los tres niveles ontológicos, patrón coherente con la lectura de la
sección~\ref{sec:dermapixel-results} sobre la complementariedad
entre el preentrenamiento auto-supervisado de PanDerm y el
contrastivo texto-imagen de DermLIP. La clasificación
\emph{zero-shot} con CLIP-L (\emph{prompts} en castellano e
inglés idénticos a la sección~\ref{sec:dermapixel-zs}) resulta
catastrófica sobre los tres niveles (AUROC L1 $0{,}559$ en
castellano y $0{,}635$ en inglés, AUROC L2 $0{,}699$ en
castellano, AUROC L3 $0{,}617$ en castellano), confirmando que los
codificadores generalistas requieren adaptación supervisada para
operar sobre material clínico dermatológico.

\section{\emph{Retrieval} $k$-NN con índice FAISS sobre
embeddings PanDerm Large}
\label{sec:dermapixel-extra-faiss}

Una alternativa al clasificador paramétrico sobre el espacio de
\emph{embeddings} congelados consiste en la recuperación basada en
similitud: dado un \emph{embedding} de consulta, se recuperan sus
$k$ vecinos más próximos del conjunto de entrenamiento bajo
similitud coseno y se predice por votación. Se construye un índice
\texttt{IndexFlatIP} de FAISS~\cite{faiss2019} sobre los
$874$ \emph{embeddings} L2-normalizados de PanDerm Large
correspondientes al conjunto de \emph{train} de la
sección~\ref{sec:dermapixel-results} y se evalúan dos variantes
(votación por mayoría y \emph{softmax} ponderado por
similitud) cruzadas con $k \in \{1, 5, 10, 20\}$. La
tabla~\ref{tab:dermapixel-extra-faiss} reporta la mejor combinación
por nivel y métrica frente a la referencia LP paramétrica.

\begin{table}[H]
\centering
\footnotesize
\setlength{\tabcolsep}{4pt}
\renewcommand{\arraystretch}{0.95}
\caption{\emph{Retrieval} $k$-NN sobre embeddings PanDerm Large
con índice FAISS (\texttt{IndexFlatIP}, similitud coseno tras
L2-normalize). Mejor combinación $k \times $ variante por nivel y
métrica. Referencia: sondeo lineal paramétrico sobre el mismo
encoder (LP, tabla~\ref{tab:dermapixel-lp}). Mejor valor por
columna y nivel en negrita.}
\label{tab:dermapixel-extra-faiss}
\begin{tabular}{llrrr}
\hline
Nivel & Método & Acc@1 & BAcc & AUROC \\
\hline
L1 & LP paramétrico (LR)              & $0{,}750$          & $\mathbf{0{,}735}$ & $\mathbf{0{,}796}$ \\
L1 & $k$-NN $k=10$ \emph{majority}     & $0{,}722$          & $0{,}686$          & $0{,}748$ \\
\hline
L2 & LP paramétrico (LR)              & $0{,}250$          & $0{,}265$          & $\mathbf{0{,}793}$ \\
L2 & $k$-NN $k=20$ \emph{majority}     & $\mathbf{0{,}278}$ & $\mathbf{0{,}265}$ & $0{,}700$ \\
\hline
L3 & LP paramétrico (LR)              & $0{,}179$          & $0{,}184$          & $\mathbf{0{,}813}$ \\
L3 & $k$-NN $k=5$ \emph{majority}      & $\mathbf{0{,}214}$ & $\mathbf{0{,}210}$ & $0{,}615$ \\
L3 & $k$-NN $k=20$ \emph{weighted}     & $0{,}214$          & $0{,}210$          & $0{,}719$ \\
\hline
\end{tabular}
\end{table}

El \emph{retrieval} $k$-NN supera al sondeo lineal paramétrico
sobre el régimen de cola larga severa L3 ($224$ clases efectivas,
$3{,}5$ muestras/clase media en \emph{train}): $+3{,}5$\,pp
Acc@1 ($0{,}179 \to 0{,}214$) y $+2{,}6$\,pp BAcc
($0{,}184 \to 0{,}210$) con $k = 5$ y votación por mayoría. Sobre
L2, $k$-NN con $k = 20$ empata en BAcc ($0{,}265$) y mejora
Acc@1 en $+2{,}8$\,pp ($0{,}250 \to 0{,}278$). Sobre L1, en
cambio, $k$-NN pierde respecto al LP paramétrico
($-4{,}9$\,pp BAcc, $-4{,}8$\,pp AUROC), comportamiento
consistente con que la cardinalidad reducida del nivel ($4$ clases)
favorece a la regresión logística sobre la votación de vecinos. La
AUROC global del $k$-NN es sistemáticamente inferior a la del LP
(diferencias entre $-4{,}8$ y $-9{,}3$\,pp según nivel),
atribuible a que la estimación de probabilidad por frecuencia
relativa de vecinos no es densa sobre el conjunto de clases. El
hallazgo es coherente con la arquitectura del módulo M4 del
prototipo de producción (\emph{retrieval} sobre Derm1M descrito en
el Anexo~\ref{cap:anexoh}), que emplea el mismo principio de
similitud coseno sobre \emph{embeddings} PanDerm Large como
soporte conceptual al diagnóstico ofrecido por las cabezas
paramétricas.

\section{\emph{Zero-shot} jerárquico: cascada L1$\to$L2$\to$L3
con propagación de errores}
\label{sec:dermapixel-extra-hierzs}

La sección~\ref{sec:dermapixel-zs} evalúa la clasificación
\emph{zero-shot} sobre los tres niveles ontológicos de forma
independiente, considerando en cada caso todas las clases
candidatas del nivel. Una estrategia alternativa, motivada por el
flujo diagnóstico clínico real, consiste en una cascada
jerárquica: predecir primero L1, restringir las candidatas L2 a las
hijas del L1 predicho, y restringir las candidatas L3 a las hijas
del L2 predicho. La jerarquía L1$\to$L2 utilizada se deriva del
dataset de \emph{train} (Inflamatoria: $25$ subcategorías L2;
Infecciosa: $5$; Tumoral: $12$; Genodermatosis: $2$). El encoder
es DermLIP v2 (el mejor \emph{zero-shot} de la
sección~\ref{sec:dermapixel-zs}), bajo \emph{prompts} en
castellano. Se contrastan tres variantes: el \emph{zero-shot}
plano de la sección~\ref{sec:dermapixel-zs} (todas las clases del
nivel como candidatas), la cascada jerárquica estricta con el L1
$\arg\max$, y una variante \emph{top-3 L1} que mantiene las tres
categorías L1 más probables como candidatas y agrega sus hijas L2.
La tabla~\ref{tab:dermapixel-extra-hierzs} compara las tres
variantes.

\begin{table}[H]
\centering
\footnotesize
\setlength{\tabcolsep}{4pt}
\renewcommand{\arraystretch}{0.95}
\caption{Comparación entre clasificación \emph{zero-shot} plana y
\emph{zero-shot jerárquico} en cascada sobre el conjunto de test de
DermapixelAI~1.0 ($N_{\mathrm{test}} = 36$ en L1/L2,
$N_{\mathrm{test}} = 28$ en L3). Encoder: DermLIP v2, prompts en
castellano. La fila ``L2 hier $|$ L1 correcto'' restringe la
evaluación a las muestras con L1 predicho correctamente. Mejor
valor por columna en negrita; la cifra condicionada se reporta
separadamente.}
\label{tab:dermapixel-extra-hierzs}
\begin{tabular}{lrrr}
\hline
Variante                                  & L1 Acc@1 & L2 Acc@1 & L3 Acc@1 \\
\hline
ZS plano (referencia, sección~\ref{sec:dermapixel-zs}) & $\mathbf{0{,}361}$ & $\mathbf{0{,}222}$ & $\mathbf{0{,}107}$ \\
ZS jerárquico (cascada estricta)           & $0{,}361$          & $0{,}139$          & $0{,}036$ \\
ZS jerárquico (top-3 L1)                   & $0{,}361$          & $0{,}194$          & $0{,}036$ \\
\hline
L2 hier $|$ L1 correcto (condicionada)     & ---                & $0{,}385$          & --- \\
\hline
\end{tabular}
\end{table}

La cascada jerárquica estricta degrada el rendimiento plano sobre
L2 en $-8{,}3$\,pp Acc@1 ($0{,}222 \to 0{,}139$) y sobre L3 en
$-7{,}1$\,pp Acc@1 ($0{,}107 \to 0{,}036$). La causa
principal es la propagación de errores desde el nivel L1: dado que
DermLIP v2 acierta L1 sólo en el $36{,}1\%$ de los casos, el
$63{,}9\%$ restante recibe en L2 un conjunto de candidatas
desalineado con la verdad de terreno, lo que invalida toda
predicción posterior independientemente del codificador. La
variante \emph{top-3 L1} amortigua parcialmente la propagación al
mantener las tres categorías L1 más probables como hipótesis activa
y recupera $0{,}194$ Acc@1 sobre L2 ($-2{,}8$\,pp respecto a
plano), sin efecto sobre L3. Condicionado a L1 correcto, la
cascada jerárquica L2 alcanza Acc@1 $= 0{,}385$, casi el doble
del plano, evidencia de que la restricción del espacio de
candidatas a las subcategorías clínicamente coherentes con la
categoría etiológica acertada incrementa la discriminación de
forma sustancial. La implicación clínica es directa: la arquitectura
jerárquica funciona si la decisión inicial sobre la familia
etiológica es robusta; con \emph{zero-shot} modesto, la cascada
acumula errores y degrada el rendimiento plano.
